\documentclass[11pt]{article}
\usepackage[margin=1in]{geometry}
\usepackage{amsmath,amssymb,booktabs,array}
\usepackage[hidelinks]{hyperref}
\usepackage{xcolor}
\setlength{\parindent}{0pt}
\setlength{\parskip}{0.55em}
\newcommand{\ve}{\text{veh-eq}}
\newcommand{\dd}{\,\mathrm{d}}
\title{S13 M2 Analytical Model Contract\\\large Source-bounded continuous mass, mandatory transfer, and finite receiving}
\author{Methodology review artifact}
\date{30 September 2026}

\begin{document}
\maketitle

\section{Purpose and claim boundary}

This contract instantiates one minimal member of the M2 family admitted by S12.
It is designed to test whether hidden lane allocation can change a useful common
lane-uniform variable-speed-limit (VSL) decision through endogenous mandatory
transfer and finite receiving. It is not a calibrated freeway model, a paper
reproduction, or an implementation specification.

The contract preserves continuous vehicle-equivalent mass, fixed lane topology,
one shared infrastructure action, exact predecision aggregate-history equality,
and complete affected-cohort time accounting. The analytical certificate that
uses this contract is a conditional existence result.

\section{Geometry, cohort, hidden state, and observer}

Consider a permanent three-to-two lane drop. For the analytical reduction, the
ending lane is explicit and the two continuing lanes are represented as one
symmetric continuing group. Let a common affected cohort have total mass
\(M>0\) \(\ve\). A hidden preparation state \(s\) is described by its
ending-lane fraction \(p_s\in(0,1)\):
\begin{equation}
 e_s(0)=p_sM,\qquad c_s(0)=(1-p_s)M.
 \label{eq:init}
\end{equation}
The compared states have different \(p_s\) but the same \(M\), route
population, geometry, future forcing, action set, and declared aggregate
predecision history.

The strong observer contains spatial aggregate density, flow, speed, cumulative
counts, and their complete declared history. It does not contain lane
decomposition. S05 supplies the conditional upstream construction in which
conservative lateral transfers cancel after aggregation, so two values
\(p_L<p_H\) can arrive at the decision boundary with identical observer history.
The decision boundary is the first entry into the fixed spatial M2 preparation
zone; M2 physics is common thereafter.

\section{Action and traversal time}

The only action is a constant lane-uniform advice speed
\(v\in[v_{\min},v_{\max}]\), with \(0<v_{\min}<v_{\max}\). The preparation
zone has common length \(L>0\) km. On the selected uncongested longitudinal
branch, advice determines zone traversal time
\begin{equation}
 \tau(v)=\frac{L}{v},\qquad
 \tau\in[\tau_{\min},\tau_{\max}]
 =\left[\frac{L}{v_{\max}},\frac{L}{v_{\min}}\right].
 \label{eq:tau}
\end{equation}
This is the full one-parameter action family. Fastest advice \(v_{\max}\) is
the endpoint \(\tau_{\min}\); it is not omitted from the comparison.

\section{Conservative mandatory transfer}

Within the preparation zone, ending-lane mass has mandatory lateral demand. On
the fixed positive active branch used for this certificate, a common transfer
rate \(\lambda>0\) (h\(^{-1}\)) gives
\begin{equation}
 \dot e_s(t)=-\lambda e_s(t),\qquad
 \dot c_s(t)=+\lambda e_s(t),\qquad 0\le t\le\tau.
 \label{eq:transfer}
\end{equation}
Thus
\begin{equation}
 e_s(\tau)=p_sM e^{-\lambda\tau},\qquad
 c_s(\tau)=M-e_s(\tau),\qquad e_s(t)+c_s(t)=M.
 \label{eq:remaining}
\end{equation}
Equation~\eqref{eq:transfer} is a deterministic continuous-mass specialization
of repeated mandatory CTM transfer on one active density/geometry branch. It is
an explicit S13 analytical assumption. It does not assert an empirically
calibrated exponential lane-change law.

Target receiving in the preparation zone is assumed nonbinding for the admitted
certificate regime. This does not remove receiving from M2: the consequential
finite-receiving interaction occurs at the choke. It isolates the mechanism by
which longer residence can complete more mandatory transfer before that
interaction.

\section{Finite receiving at the choke}

At the choke, same-lane mass consumes one unit of receiving work per unit mass.
Mass that still must merge consumes \(\alpha>1\) units, following the
gap-acceptance receiving-work interpretation in Carey, Balijepalli, and Watling.
Let \(\mu>0\) be available receiving work per hour. The complete cohort's
receiving workload is
\begin{equation}
 W_s(\tau)=c_s(\tau)+\alpha e_s(\tau)
 =M+(\alpha-1)p_sM e^{-\lambda\tau}.
 \label{eq:work}
\end{equation}
Neutral proportional allocation shares binding receiving between effective
same-lane and gap-weighted lateral demand. For the homogeneous fluid batch used
here, workload is served at rate \(\mu\), so clearance duration is
\begin{equation}
 T_s(\tau)=\frac{W_s(\tau)}{\mu}.
 \label{eq:clear}
\end{equation}
The service rule is conservative and work-conserving in receiving-work units.
It does not create or delete vehicle mass. The gap factor changes time/space
consumption, not mass.

\section{Complete affected-cohort objective}

All \(M\) mass begins at the preparation-zone entrance at the decision time,
traverses for \(\tau\), and then clears the choke under homogeneous proportional
fluid service. Departures are linear over \(T_s\), so mean post-arrival time in
system is \(T_s/2\). The per-mass complete time score is
\begin{align}
 J_s(\tau)
 &=\tau+\frac{T_s(\tau)}{2} \\
 &=\tau+\frac{M}{2\mu}
   +Kp_s e^{-\lambda\tau},
 &K&:=\frac{(\alpha-1)M}{2\mu}>0.
 \label{eq:objective}
\end{align}
This score includes the full action-induced traversal delay and the time until
the complete cohort clears. No unfinished demand or upstream waiting is omitted.
A common action-independent prefix cost or downstream travel time can be added
without changing the comparison.

\section{Parameter class and active regime}

The global parameter domain is
\begin{equation}
 M,L,\mu,\lambda>0,\quad \alpha>1,\quad
 0<v_{\min}<v_{\max},\quad 0<p_L<p_H<1.
 \label{eq:domain}
\end{equation}
S13 does not select one calibrated point. The action-split certificate requires
the open inequalities
\begin{equation}
 K\lambda p_L<e^{\lambda\tau_{\min}}
 <K\lambda p_H<e^{\lambda\tau_{\max}}.
 \label{eq:open}
\end{equation}
Strict inequalities make this a nonempty open regime whenever the common
parameters permit ending-lane fractions on both sides of the threshold. They
are declared before evaluating an action ranking. Equality surfaces are excluded
active-set boundaries, not silently rounded cases.

\section{Aggregate limits and falsification}

The hidden-state contribution disappears when any of the following holds:
\begin{itemize}
\item \(\alpha\to1\): lateral and same-lane traffic consume equal receiving work;
\item \(\lambda\to0\): advice no longer changes remaining mandatory transfer;
\item \(p_L=p_H\): the hidden preparations coincide; or
\item receiving never binds: the workload term no longer delays clearance.
\end{itemize}
In these limits, fastest advice is common best and S05's obstruction is recovered.

The M2 claim is falsified if the exact observer collision cannot reach the zone
entrance, the shared speed cannot change residence/remaining transfer, complete
accounting removes the action split, or an observer-only deterministic model can
condition on information that distinguishes the preparations at decision time.

\section{Differentiability boundary}

Inside the fixed active regime, equations~\eqref{eq:tau}--\eqref{eq:objective}
are smooth in continuous state and action. Lane topology and the ending-lane
label are fixed metadata. Action endpoints, receiving becoming nonbinding,
priority-rule changes, zero mass, and CTM demand/supply branch switches are
nonsmooth or active-set boundaries. S13 makes no smoothing claim across them.

\section{Source and adaptation statement}

Per-lane conservation and mandatory transfer are grounded in
\href{https://doi.org/10.1080/21680566.2019.1700846}{Nagalur Subraveti,
Knoop, and van Arem}. Finite receiving, proportional priority, and the
lane-change gap factor are grounded in
\href{https://doi.org/10.1007/s11067-013-9193-7}{Carey, Balijepalli, and
Watling}. Common VSL at a permanent lane drop is grounded in
\href{https://doi.org/10.1080/23249935.2022.2033878}{Zhang et al.}
The constant-rate active branch, batch handoff, and certificate are S13's
explicit conditional analytical adaptations.

\end{document}
